\section{Контекст}
\label{sec:background}

\subsection{Память со структурированным состоянием}

\SkillState{}~\citep{skillstate2026} поддерживает устойчивое структурированное
состояние $\Sigma_t$ и собирает каждый промпт как
\[
  \text{промпт}_t \;=\; [\,\text{спецификация задачи}\,] \;+\; [\,\Sigma_t\,] \;+\;
  [\,o_t\,],
\]
где $o_t$ — последнее наблюдение, а история взаимодействия отбрасывается.
Модель выдаёт рассуждение, патч состояния и действие за одну генерацию,
$(R_t, \Delta\Sigma_t, a_t)$, как один объект JSON ровно с двумя ключами
\texttt{state\_patch} и \texttt{action}\origin. Рантайм проверяет каждый
патч по схеме и отклоняет недействительные с циклом отката и повтора. Схемы
пишутся один раз на домен; на 100 экземплярах InterCode CTF повторно
используется одна пятиполевая схема\origin.

Приводимое доказательство — это развёртка по горизонту на процедурном
генераторе: размер промпта остаётся около \num{1800} символов от $T{=}10$ до
$T{=}200$, тогда как базовые линии растут до \num{48000}, а качество
сохраняется или улучшается\origin. Собственная таксономия ошибок статьи для
моделей с открытыми весами приписывает \SI{68}{\percent} отказов
преждевременной перезаписи или удалению состояния\origin — цифра, к которой
мы вернёмся в \S\ref{sec:composition}, потому что это аргумент в пользу
вывода части состояния из-под контроля модели.

\subsection{Детерминированное состояние выполнения}

\Ledger{}~\citep{ledger2026} занимает противоположный угол дизайна:
поддерживаемая рантаймом запись, механически выводимая из завершённых
взаимодействий: записи наблюдений, счётчики изменений по файлам и глобальные,
нормализованные записи команд с их исходами. Он не добавляет обращений к
модели и не изменяет агента. Его два механизма аккуратно разделяются в его
собственной абляции: \emph{inform} (отрисовка записи в конце входа, где она
не может нарушить дополняемый префикс) даёт основную часть выигрыша по
эффективности, тогда как \emph{govern} (перехват предложенной команды и
повторное использование более раннего результата) даёт основную часть
выигрыша по разрешению\origin.

Две системы дополняют, а не соревнуются друг с другом: \SkillState{}
удерживает семантическое состояние в модели, тогда как \Ledger{} удерживает
объективное состояние в рантайме. Наше исследование композиции
(\S\ref{sec:composition}) варьирует именно эту границу, и мы намеренно
исключаем из него \emph{govern}: govern меняет то, что агент \emph{делает},
так что его включение примешало бы крупный эффект к измерению меньшего.

\subsection{Проверяемые предположения}

Сформулированы как проверяемые утверждения, а не как критика:

\begin{description}[style=nextline]
\item[A1 (доступность схемы)] Адекватная доменная схема существует до
  выполнения. Оригинал не даёт метода проектирования схемы и приводит
  дословный список полей лишь для одного из четырёх своих доменов\origin.
\item[A2 (положительный эффект с первого хода)] Поддержание состояния в
  худшем случае нейтрально с первого хода. Кратчайший приводимый горизонт —
  $T{=}10$, что уже длиннее значительной доли реальных программных задач.
\item[A3 (достаточная статистика)] Ничто отброшенное никогда больше не
  понадобится. Оригинал называет это ограничением\origin{}, но не измеряет
  его. Именно это предположение фальсифицирует данная статья.
\item[M (валюта измерения)] Совокупное число токенов — валидный показатель
  для стоимости. Полнотекстовый поиск по оригиналу даёт ноль вхождений слов
  cache, billing или иной денежной терминологии\origin.
\end{description}

A1 и A3 — свойства метода. M — свойство того, \emph{как} метод оценивался, и
именно оно затрудняет видимость двух других: метрика, считающая логические
токены, не может отразить ни скидку кэша, которая вознаграждает дополняемый
префикс, ни вызовы инструментов, потраченные на повторное приобретение того,
что промпт больше не несёт.

\subsection{Ось дизайна, используемая в этой статье}

Вместо сравнения именованных архитектур мы фиксируем механизм и варьируем
одно целое число: количество $\tailk$ самых последних дословных обменов,
сохраняемых рядом с блоком состояния.
\[
  \text{промпт}_t \;=\; [\,\text{спецификация задачи}\,] \;+\;
  [\,\text{последние } \tailk \text{ обменов}\,] \;+\; [\,\Sigma_t\,]
\]
При $\tailk{=}0$ это в точности \SkillState{} в опубликованном виде. По мере
роста $\tailk$ конфигурация интерполирует в сторону агента с транскриптом.
Контрольная группа не является точкой на этой оси (у неё вообще нет блока
состояния) и всегда приводится отдельно.

Это даёт три вещи, которые не может дать трёхрукое сравнение. Оно устраняет
произвольную константу, которую незаметно внесла бы промежуточная
«гибридная» рука; оно превращает вопрос из «какая архитектура» в «где на оси
лежит оптимум»; и оно делает форму зависимости наблюдаемой, что в итоге и
локализует ущерб (\S\ref{sec:capability}).
